\documentclass[../../../include/open-logic-section]{subfiles}

\begin{document}

\olfileid{sth}{z}{story}
\olsection{سیٹ بنن دا بیان ہور تفصیل نال}

\olref[story][approach]{sec} وچ اسیں شون فیلڈ دا سیٹ بنن دی کارروائی
بارے قول دتا سی۔ ہُن اسیں کجھ ہور اصول لکھنے چاہندے آں تاں جو ایہ بیان
ذرا ہور ٹھیک ہو جاوے۔ اصول ایہ نیں:
\begin{enumerate}
\item[] \stageshier۔ ہر سیٹ کسے مرحلے اُتے بندا اے۔
\item[] \stagesord۔ مرحلے ترتیب وچ نیں، کجھ ہورناں توں \emph{پہلاں}
آؤندے نیں۔\footnote{اسیں اصل وچ بناں آکھے ای منّاں گے کہ مرحلے
\emph{حُسنِ ترتیب} وچ نیں۔ ایہدا مطلب \olref[ordinals][]{chap} وچ
سمجھایا اے۔ ایہ اک وڈا مفروضہ اے۔ اصل وچ \citet{Scott1974} دی اک
بڑی سُجھ والی ترکیب نال ایہ مفروضہ پہلاں \emph{چھڈیا} تے فیر
\emph{اخذ کیتا} جا سکدا اے۔ (ایس نال ایہ وی سمجھ آوے گی کہ سانوں
پہلا مرحلہ موجود کیوں سمجھنا چاہیدا اے۔) ایہ گل ایتھے نہیں کھول سکدے۔
ہور لئی \citet{ButtonLT1} ویکھو۔}
\item[] \stagesacc۔ کسے وی مرحلے $S$ لئی، تے مرحلے~$S$ توں \emph{پہلاں} بنے سیٹاں دی کسے وی چُناؤ لئی، مرحلے~$S$ اُتے اک
سیٹ بندا اے جیہدے رکن عین اوہی سیٹ نیں۔ مرحلے~$S$ اُتے ہور کجھ نہیں بندا۔
\end{enumerate}
ایہ غیر رسمی اصول نیں، پر اسیں ایہناں نال زرمَیلو دے سیٹ نظریے دے کئی
مسلّمیاں دا جواز دے سکاں گے۔

(اک گل دا دھیان دئیے۔ اسیں کجھ بالکل معیاری مسلّمے، بالکل معیاری ناںواں
نال، پیش کراں گے۔ پر ہُنے دتے ٹیرھے حرفاں والے اصولاں دے علمی لکھت وچ
کوئی خاص ناں نہیں۔ اسیں بس اوہ ناں ورتدے آں جیہڑے ساڈے خیال وچ کم آؤن گے۔)

\end{document}
